% !TEX root = ../Thesis.tex

\chapter{Introduzione}
L'esigenza di una comunicazione riservata, autentica e integra costituisce un pilastro fondamentale dell'architettura del mondo digitale contemporaneo. Tradizionalmente, la sicurezza delle informazioni è stata affidata alla crittografia e a protocolli standard come TLS/SSL, progettati per rendere il contenuto dei messaggi inaccessibile a chiunque non possieda le chiavi di decifratura.

Tuttavia, l'aumento delle misure di sorveglianza di massa e di censura capillare ha messo in luce una vulnerabilità critica intrinseca alla crittografia: la sua visibilità. Sebbene il contenuto rimanga segreto, l'elevata entropia del traffico cifrato lo rende facilmente distinguibile dalle comunicazioni in chiaro. Questo permette a regimi autoritari o a sofisticati meccanismi di \emph{firewalling} di identificare, bloccare o degradare selettivamente le comunicazioni protette, rendendo di fatto impossibile lo scambio di informazioni in contesti ostili.

In scenari di censura estrema emerge dunque la necessità di ricorrere alla steganografia. Questa disciplina si occupa dell'occultamento di informazioni segrete all'interno di oggetti di copertura apparentemente innocui (come testi, immagini o audio), detti \emph{cover object}. A differenza della crittografia, il cui successo dipende dalla difficoltà matematica di decifrare il messaggio, la sicurezza steganografica dipende dall'indistinguibilità statistica: un osservatore esterno non deve essere in grado di discriminare tra un oggetto contenente un messaggio segreto (\emph{stego-object}) e uno privo di esso.

Le metodologie steganografiche classiche, come la modifica dei bit meno significativi (LSB) nelle immagini, si sono rivelate vulnerabili alla moderna crittoanalisi statistica. Per superare queste limitazioni e garantire una sicurezza dimostrabile, la ricerca recente ha operato un cambio di paradigma, orientandosi verso l'utilizzo di modelli generativi.

In questo contesto si colloca \textbf{Meteor}, un algoritmo di steganografia a chiave simmetrica che sfrutta le capacità dei modelli generativi (originariamente modelli di linguaggio come GPT-2) per creare canali di comunicazione occulti. Meteor non modifica un contenuto esistente, ma guida il processo di generazione stesso: utilizza la distribuzione di probabilità del modello per campionare i token successivi in base al messaggio segreto cifrato. Questo garantisce che lo \emph{stego-object} risultante sia statisticamente coerente con l'output naturale del modello, rendendo la presenza del messaggio estremamente difficile da rilevare.

L'obiettivo di questa tesi è estendere il principio di sicurezza steganografica di Meteor all'ambito delle immagini ad alta risoluzione.
Per raggiungere tale scopo, viene sfruttata l'efficienza dell'architettura \textbf{VQ-GAN} (Vector Quantized Generative Adversarial Network). La scelta di questa architettura è strategica: la VQ-GAN permette di modellare immagini complesse attraverso un vocabolario discreto di "token visivi" (codebook), consentendo l'applicazione di un Transformer autoregressivo per la generazione.
Il lavoro svolto mira a dimostrare come sia possibile incorporare messaggi cifrati durante la fase di generazione dell'immagine da parte della VQ-GAN, ottenendo risultati visivamente realistici e robusti contro i tentativi di rilevamento.

Le implicazioni di questo approccio sono significative per lo sviluppo di strumenti \emph{censorship-resistant}: abilitare lo scambio di informazioni nascoste in immagini di alta qualità offre un nuovo livello di protezione per la privacy degli utenti e la libertà di espressione in ambienti digitali sorvegliati.

\medskip


Il presente elaborato è organizzato come segue:

\begin{itemize}
	\item Il \textbf{Capitolo 2} fornisce il background teorico necessario, analizzando lo stato dell'arte della steganografia e i principi di funzionamento dei modelli generativi, con particolare focus su GAN e Transformer.
	
	\item Il \textbf{Capitolo 3} descrive in dettaglio l'algoritmo Meteor e l'architettura VQ-GAN, ponendo le basi teoriche per la loro integrazione proposta in questa tesi.
	
	\item Il \textbf{Capitolo 4} presenta l'architettura del sistema implementato, illustrando le modifiche apportate al meccanismo di campionamento per permettere l'embedding del messaggio nel flusso di token visivi.
	
	\item Il \textbf{Capitolo 5} espone i risultati sperimentali, valutando la qualità visiva delle immagini generate (tramite metriche come FID) e la capacità di steganografica (payload e robustezza).
	
	\item Infine, il \textbf{Capitolo 6} trae le conclusioni del lavoro svolto e delinea possibili sviluppi futuri per migliorare l'efficienza e la sicurezza del sistema proposto.
\end{itemize}